\section{Coding Strategy}

The previous section defined the case-level research design. This section
describes how the main outcome is measured. The coding strategy begins from a
simple premise: official exit language is evidence about regulatory status, not
by itself evidence that a government financing function disappeared. For each
platform, I therefore code both the formal event and the post-event location of
the relevant function.

\subsection{Source hierarchy}

The source hierarchy begins with government documents, finance-bureau notices,
state-owned asset supervision documents, exchange disclosures, and company
announcements when they are available. These sources are most useful for
identifying official exit, market-oriented transformation, asset transfer, debt
replacement, restructuring, and liquidation. Bond prospectuses, credit rating
reports, legal opinions, annual reports, and audited financial statements are
often more informative about continuing public-project functions, repayment
arrangements, fiscal support, issuer guarantees, related-party transfers, and
changes in the role of the company after exit. Credible financial news and
company profiles serve as supplementary sources. These materials can help
locate events and identify entities, but they do not by themselves determine
the final label.

This hierarchy matters because the most authoritative source for the formal
event is not always the most informative source for the post-event function. A
government notice may state that a platform exited a list, while a later
prospectus may show that the same issuer still carries out infrastructure
investment, entrusted construction, land preparation, or project repayment on
behalf of the local state. Conversely, a prospectus may contain
no-government-financing language, while the underlying fiscal-substitution
arrangement appears only in a debt-replacement announcement or a municipal
budget document. The coding rule therefore treats each source as evidence for
a specific component of the case rather than as a global statement about the
whole outcome.

\subsection{Condensed codebook}

The full coding protocol is maintained in the project codebook. The condensed
rule used in the paper begins by identifying the baseline platform function.
Relevant functions include infrastructure investment and financing,
land development, government-entrusted construction, policy-driven housing or
urban renewal, public utility construction, and repayment arrangements tied to
local fiscal resources. The next step is to identify the formal event, such as
exit from a financing-platform list, a no-government-financing statement,
market-oriented transformation, merger, asset transfer, debt replacement, debt
restructuring, deregistration, liquidation, or dissolution. The final coding
decision turns on the post-event location of the function. The same issuer may
continue the function, the function may move into formal fiscal channels, the
function may move to another local state-owned entity, or the original legal
entity may terminate.

A case enters the final label file only if it satisfies two requirements. It
must contain a formal event, and it must contain post-event evidence about the
relevant function. Cases with only a corporate profile, only a news summary,
only an official slogan of market-oriented transformation, or only
post-event public project evidence without a formal event are retained as
candidates rather than final labels. Purely commercial SOEs, central SOEs, and
provincial entities without a city-platform role are excluded unless the source
packet shows a platform-like local financing function.

\begin{table}[htbp]
\centering
\footnotesize
\setlength{\tabcolsep}{4pt}
\begin{tabular}{@{}>{\raggedright\arraybackslash}p{0.18\textwidth}>{\raggedright\arraybackslash}p{0.36\textwidth}>{\raggedright\arraybackslash}p{0.34\textwidth}@{}}
\toprule
Label & Positive rule & Exclusion rule \\
\midrule
Substantive exit & The formal event is followed by evidence that the issuer no
longer performs government financing or quasi-fiscal financing functions. Former
public projects may continue only if the financing has moved into budgetary
channels, special-purpose bonds, or fee-based project management. & Do not code
substantive exit from no-government-financing language alone. Continuing fiscal
repayment, entrusted construction financing, land preparation financing, or
public project financing rules it out. \\
Nominal exit & The same issuer has formal exit or compliance language but
continues the relevant platform function after the event. Evidence may include
infrastructure financing, land preparation, shantytown redevelopment, fiscal
receivables, government purchase service, project repurchase, or fiscal
subsidies. & Do not code nominal exit when the function clearly moved to a
different local state-owned entity, or when post-event sources show only
fee-based project management funded through formal fiscal channels. \\
Functional transfer & The original issuer exits, reduces its role, merges, or
is reorganized, while the relevant financing, construction, debt-management, or
public project function moves to another local state-owned entity, subsidiary,
merged group, or newly created platform. & Do not code functional transfer if
the same issuer continues the function. Do not code it when the function is
absorbed into formal fiscal channels rather than another platform-like entity.
\\
Liquidation & The firm is dissolved, deregistered, liquidated, bankrupted, or
otherwise terminated, with no evidence that the same legal entity continues
platform functions. & If a successor local SOE takes over the fiscal function,
code functional transfer and record the liquidation as part of the transfer
process. \\
\bottomrule
\end{tabular}
\caption{Condensed codebook for LGFV exit types}
\label{tab:condensed-codebook}
\end{table}

Several edge cases matter for the classification. Fee-based project management
after exit is not automatically nominal exit. If the project is funded through
formal fiscal channels and the issuer receives only a management fee, the case
is closer to substantive exit. If the issuer carries project financing,
government repayment exposure, or fiscal receivables, the case is closer to
nominal exit. Debt replacement also has two meanings. It supports substantive
exit only when legacy obligations move into formal fiscal management and no
new quasi-fiscal function is visible. If debt replacement coexists with
continued platform financing, the case remains nominal exit. Mergers and state
capital operation companies require special care. If a public financing
function is preserved inside a reorganized local SOE system, the case is coded
as functional transfer rather than substantive exit.

Each final label is accompanied by three scores. The confidence score records
whether the final label is high, medium, or low confidence. High confidence
requires a directly documented formal event, clear post-event function
evidence or rich negative evidence, and no plausible alternative label. Medium
confidence permits a partial source packet or a credible but weaker
alternative interpretation. Low confidence cases do not enter the final label
file. The source coverage score ranges from 0 to 4 and measures whether the
case has original or near-original sources for both the formal event and the
post-event function. The continued-function evidence score also ranges from 0
to 4 and measures the strength of evidence that platform-like functions
continued after the event. A high continued-function score does not
mechanically imply nominal exit because the continued function may be located
in a successor entity rather than the same issuer.

Table~\ref{tab:pilot_cases} summarizes five illustrative working-reference
cases. The full file now contains ninety-four human-checked reference labels; this
table is not an empirical result or an exhaustive case list. Its purpose is to
show how the same coding rule operates across different documentary patterns.

\begin{table}[htbp]
\centering
\footnotesize
\setlength{\tabcolsep}{4pt}
\begin{tabular}{@{}>{\raggedright\arraybackslash}p{0.18\textwidth}>{\raggedright\arraybackslash}p{0.22\textwidth}>{\raggedright\arraybackslash}p{0.34\textwidth}>{\raggedright\arraybackslash}p{0.13\textwidth}@{}}
\toprule
Case & Source basis & Key evidence & Final label \\
\midrule
Xi'an Hi-tech Holding & 2020 and 2026 bond prospectuses, rating report, and
supplementary financial news & The issuer reports a 2012 exit from the
financing-platform list, while later documents continue to describe
infrastructure investment, BT, entrusted-construction, and government-linked
project repayment. & Nominal exit \\
Chengdu City Construction Investment Management & 2023 bond prospectus and
Shanghai Clearing disclosure page & The issuer states that it does not
undertake government financing functions after 2015, while the same prospectus
describes continuing infrastructure investment, fiscal payment arrangements,
self-funded and financed project construction, and local debt replacement. &
Nominal exit \\
Guangzhou City Construction Investment & 2018 tracking rating report, 2018
debt-replacement announcement, and 2026 bond prospectus & The company reports
that it stopped new government-project investment and financing after 2014,
shifted to fiscal-funded entrusted management, and had part of a 2014 MTN
included in the fiscal budget for debt replacement. & Substantive exit \\
Hangzhou City Investment & 2024 bond prospectus and issuer rating report &
Qiantou Group was folded into Hangzhou City Investment, Hangzhou City
Investment succeeded Qiantou debt instruments, and Anju Group was created as a
housing-security investment, construction, and operation platform. &
Functional transfer \\
Zhuzhou State-owned Assets Investment Holding & 2026 bond prospectus and legal
opinion & The issuer reports a March 2013 exit from local government financing
platform management, while the same source set describes land preparation,
major project construction, government-entrusted work, fiscal settlement,
subsidies, and local government support. & Nominal exit \\
\bottomrule
\end{tabular}
\caption{Illustrative coding examples}
\label{tab:pilot_cases}
\end{table}

These cases illustrate distinctions that are central to the measurement
strategy. Formal compliance language must be separated from functional
evidence. A prospectus may state that new debt does not constitute local
government debt, while other sections describe continuing public project
obligations. Chengdu makes this distinction especially clear because the
no-government-financing statement appears in the same document as evidence of
self-funded and financed infrastructure construction with fiscal payment
arrangements. Continuing public project management is also not always evidence
of nominal exit. In Guangzhou, the relevant distinction is that government
projects were funded through municipal fiscal funds and the company received a
management fee, rather than using the platform as an off-budget financing arm.
Functional transfer, finally, is not merely a name change. In Hangzhou, the
classification rests on equity transfer, debt succession, and the creation of a
specialized housing-security group within a reorganized municipal platform
system.

\subsection{LLM-assisted screening}

Large language models are used as a screening device, not as final outcome
coders. Each prompt applies the same codebook to a source excerpt and returns a
proposed label, confidence level, evidence quotation, rationale, and most
plausible alternative. It must return \texttt{unclear} when the packet lacks
either a formal event or evidence about the post-event location of the
function. This requirement reduces the temptation to infer institutional
change from generic compliance language.

The expanded screening file contains 361 disclosure-level candidates. Fifteen
lack a usable source packet. Among the remaining 346 rows, 94 belong to the
working-reference set, 203 receive one-sided LLM surrogate labels for
nominal exit, 44 contain no directly documented formal exit or compliance
event, and five are source-packet-reviewed boundary cases. The 203 surrogate rows
correspond to 158 issuers because some platforms issue repeated disclosures.
Sixty-one of these issuers also appear in the reference set, while 97 do not.
Rows without a documented formal event remain screening observations rather
than negative exit outcomes.

The current surrogate rule is intentionally narrow. It identifies nominal-exit
candidates only when a document combines direct no-government-financing,
no-new-government-debt, or no-hidden-debt language with evidence that the same
issuer continues infrastructure financing, entrusted construction, land
development, fiscal support, or another public-project function. The screen
does not attempt to infer substantive exit, functional transfer, or liquidation
from silence. It therefore expands source review for a common form of formal
compliance, but it is not a four-category outcome measure.

\subsection{Measurement audit}

The working-reference file contains 94 cases assigned through Codex review of
the retained source packets on behalf of the project author under the frozen
codebook. On 10 September 2026, the author reported that all 94 labels had
subsequently been checked by a human and required no revisions. The date
records the author's confirmation, not the date on which checking occurred.
The original Codex labels are preserved, and a case-level confirmation register
links the human-check report to that exact snapshot. The
labels comprise two substantive exits, 82 nominal exits, ten functional
transfers, and no liquidations. Every
evidence identifier in the file resolves to a source inventory record.
Ninety-three cases have complete local copies of all cited documents, all 94
have recovery URLs for every cited document, and 62 have an exact case-level
evidence memo. These figures establish source traceability. They do not by
themselves establish label accuracy or intercoder reliability.

The 61 issuers that overlap the LLM screen and reference set receive a nominal
label under both procedures. This 61-of-61 agreement is useful as a
selected-overlap consistency check: the screen can recover documentary patterns
that the codebook treats as nominal exit. It is not an estimate of population
precision because the overlap was assembled after positive screening rather
than drawn under known inclusion probabilities. No issuer screened as lacking
a direct formal event currently overlaps the reference labels. The existing
data therefore cannot estimate recall, four-category calibration, or the error
rate among screen-negative cases. No independent pre-adjudication coding
forms have been supplied. Reviewer identity, the actual review date, and
blinding are not documented by the author's report, so the paper does not
report an intercoder statistic or treat the reported absence of revisions as
a measured accuracy rate.

These boundaries also limit the use of surrogate-label estimators. The
framework developed by \citet{egami2023dsl} treats predicted labels as noisy
surrogates and uses a probability validation sample with human outcomes to
adjust downstream estimates. That logic is appropriate for the intended design,
but the present overlap does not supply the required validation probabilities
or independently recorded validation-sample outcomes. Human checking of the
existing reference file does not remedy that selection problem. The paper
therefore reports the LLM output as a
screening flow and does not present a doubly robust or validation-adjusted
effect estimate.

\subsection{Validation frame and collection}

The validation frame begins with 133 distinct issuers outside the reference
overlap: 97 screen-positive issuers and 36 whose reviewed disclosures lack a
direct formal event. Source review resolves the 98 inherited scope decisions
and uniquely assigns 87 of the 88 inherited geography gaps. One already
ineligible commercial issuer retains multiple locations under the approved
geography rule. After scope exclusions, 67 issuers remain, comprising 57
screen-positive issuers and ten without a directly documented formal event.
The current citations are verified against 159 sources and 488 evidence
locators. This completes frame construction, not outcome validation.

The existing allocation proposes reviewing all 67 issuers across 24 strata
defined by screening status, source coverage, historical-capacity availability,
debt-data availability, and administrative level. Within this assembled frame,
the proposed inclusion probability is one. Because each stratum is taken in
full, the current allocation does not implement differential oversampling.
The original proposal and subsequent freeze decisions are retained separately;
a proposed inclusion probability is not a record of completed human review.
The frame does not represent Chinese LGFVs generally or all issuers in the
original disclosure pool.

Collection separates platform scope from exit-case eligibility. A reviewer
must first establish the formal event and post-event function before assigning
one of the four exit types. A packet without sufficient event evidence remains
unresolved or outside the exit-case outcome analysis; it is not coded as a
negative institutional-change outcome. The ten cases without direct event
language are retained so that human review can examine what the screen missed,
rather than making the screening decision determine its own validation.

The collection materials separate a blinded source sheet from a coordinator
file containing the screening result and design variables. Human decisions,
event and function evidence, uncertainty, and adjudication are recorded by
stable issuer and case identifiers. Existing human-checked labels can be reused
only when the issuer and event correspond to the same case. A shared city or
similar platform name is insufficient. Agreement and error rates will be
reported only for their observed denominators, with unresolved cases and
missing responses retained in the collection flow.
